\section{Билет 6. Завершающий период Гражданской войны (1919--1922 гг.) и её итоги}

\subsection{Основные события Гражданской войны на восточном и южном фронтах в 1919--1920 гг.}

\subsubsection{Юго"=западный фронт}

В конце 1918 года на Украине возникает Директория Украинской НР во главе с Петлюрой. Она оказывает сопротивление как большевикам, так и армии Деникина.
Армия Директории (Украинская народная армия) состояла из 30 тысяч человек.

Советско"=украинский фронт предпринимает наступление в январе--феврале 1919 года. Он занимает Харьков, Киев и продвигается в сторону Львова.

В это же время войска интервентов постепенно оставляют Херсон, Николаев и отступают к Одессе. В этих войсках под влиянием советской пропаганды начинаются волнения.

В конце марта 1919 года состоялась Парижская мирная конференция, на которой союзники стран Антанты принимают решение об эвакуации войск с территории России.

\subsection{Восточный фронт}
Масштабные боевые действия в марте 1919 года. Колчак предпринимает массированные наступления нескольких армий. Он отбивает Уфу и другие города, движется к Волге. Он требует от войск отбросить войска Красной армии за Волгу, попытаться разгромить на подходах к ней и захватить важнейшие переправы. Это вызывает ответные действия советского правительства. Ленин призывает направить все силы на разгром армий Колчака. 

Восточным фронтом стал руководить М.~В.~Фрунзе (после назначения Каменева на пост главнокомандующего). Происходит постепенное усиление Восточного фронта, начинается контрнаступление. Красная армия занимает Бугуруслан, разбивает корпус генерала Капеля (?). Большую роль сыграла армия Чапаева. Его дивизия успешно воевала с казаками до 5 сентября 1919 года, когда он был убит.

Однако наступление Колчака захлебнулось уже к апрелю 1919 года. Красная армия Восточного фронта многочисленна и не оставляет Колчаку практически никаких шансов. После взятия Ижевска окончательно стратегическая инициатива переходит к Восточному фронту.
В январе 1920 года армия Колчака была разгромлена. Части Красной армии выходят через юг Сибири к Байкалу. Но дальше они не проходят, так как к востоку от Байкала находятся армии Японии и США. В апреле 1920 года создаётся Дальневосточная республика.

В Средней Азии создаётся туркестанский фронт во главе с Фрунзе. Он нанёс окончательное поражение белым. Воспользовавшись междуусобной войной туркестанских племен, создается Хорезмская НР, а также Бухарская НР.

\subsubsection{Снова Южный фронт}

Войска Южного фронта, получившие подкрепление к 1919 году, переходят в наступление против Деникина.
24 января после оттеснения было отправлено циркулярное письмо войскам на Дону с требованием провести массовый террор против казаков, принимавших участие в борьбе с советской властью.
Ответом на такие действие стало восстание казаков в марте 1919 года, серьезно затруднившее наступление войск Южного фронта.

Деникин, готовясь к наступлению войск Южного фронта, пытается добиться централизации управления всеми антисоветскими силами. Он объединяется с донской армией Краснова и Черноморским флотом. Так появляются Вооруженные силы Юга России (ВСЮР) под командованием Деникина.

В июле 1919 года, пока силы РККА сосредоточены на разгроме Колчака, Деникин организовывает мощное наступление. ВСЮР должны двигаться через Курск и Воронеж с выходом к Москве. В июле 1919 года на пленуме РКП(б) зачитывается письмо Ленина, которое начинается как <<Все на борьбу с Деникиным>>. Южный фронт получает пополнение в виде состава, патронов и снарядов. Тем временем, ВСЮР к сентябрю захватили Курск и Воронеж, подошли к Орлу. Только тут красной армии удалось остановить армию Деникина и перейти в контрнаступление. В январе 1920 г. Вооруженные силы Юга России были окончательно разгромлены. Деникин 4 апреля 1920 года передал командование генералу Врангелю, который начал укрепляться в Крыму.
К 1920 году были разбиты основные противобольшевистские силы, что позволило укрепиться власти большевиков.

\subsection{Советско"=польская война.}
25 апреля 1920 г. польская армия вторглась на территорию советской Украины, 6 мая овладела Киевом и к середине мая вышла на левый берег Днепра. Действиями главы польского государства маршала Ю.~Пилсудского руководила идея создания <<Великой Польши>>, включающей не только польские, но и украинские, белорусские, литовские территории. Откровенная польская агрессия вызвала подъем патриотических чувств в России. Во второй половине мая войска двух советских фронтов "--- Западного (командующий М.~Н.~Тухачевский) и Юго"=Западного (А.~И.~Егоров) "--- перешли в наступление. Им удалось отбить Киев.

Далее было принято решение идти на Варшаву, надеясь на лояльность польского пролетариата. Надежды не оправдались, поляки встретили РККА враждебно и перешли в контрнаступление, выйдя на территорию Украины. Советское командование было вынуждено заключить перемирие в октябре 1920 г., а в марте 1921 г. был заключён Рижский мир, по которому Польше отходила часть украинских и белорусских земель.

\subsection{Правительство юга России и П.~Н.~Врангель}

После заключения мира, Красная армия сосредоточила силы на разгроме Врангеля в Крыму. К началу ноября 1921 года проведена операция, завершившаяся занятием Крымского полуострова и разгромом группировок Врангеля, после чего остатки белогвардейских сил эвакуировались в Турцию.
Ещё один центр войны "--- Закавказье. В результате боевых действий Красной армии создаются республики в Армении, Грузии и Азербайджане.

\subsection{Наступление Юденича на Петроград}
Летом и осенью 1919 года происходит наступление армии Юденича на Петроград, которое, под руководством самого Троцкого, удаётся отбить. К 1920 году удаётся отбросить Юденича к границам Эстонии, а в феврале РККА вступает на её территорию.

\subsection{Советская власть и крестьянство. Рабочие в годы <<военного коммунизма>>. Сопротивление советской власти на завершающем этапе войны: Кронштадтское, Тамбовское и Западносибирское восстания}

На рубеже 1918--1919 были отменены комбеды, что уменьшило сопротивление крестьянства, однако продолжались топливный и продовольственный кризисы. В городах вводится нормированное снабжение продовольствием и товаром (карточная система). Ленин выдвигает лозунг <<Не смей командовать середняком!>>. В соответствии с этой новой политикой ВЦИК и Совнарком принимают решение об освобождении от натурального налога в 36 губерниях.

Сохранение экстренных мер военно"=коммунистического характера вызвало недовольства. Крестьяне были недовольны сохранившейся продразвёрсткой. К ним присоединяются рабочие. Начинаются политические восстания. Особенно крупное "--- Кронштадтское в марте 1921 года.

\subsection{Милитаризация промышленности}

Снабжение армии осуществлялось за счет национализации производств и трудовой повинности, за счёт продовольственных отрядов, изымавших излишки у кулаков, вводится карточная система. Действует система продразвёрстки (обязательства, возлагаемые на крестьянство, поставлять продукцию государству, фактически не ограничивался размер продразвёрстки).

Трудовая повинность была первоначально только у буржуазии, но с октября 1918 г. все лица от 16 до 50 лет ставились на учёт в отделы распределения рабочей силы.

\subsection{Кризис военно"=коммунистической системы}
Экономическое положение к весне 1921 года невероятно тяжелое. Весной 1921 года начинается голод, особенно тяжелый в Поволжье и южном Приуралье. От него пострадало 20\% населения, и 5 млн человек погибло. Советское правительство обращается за помощью к другим странам. Наибольшую помощь внесла американская компания ARA.

\subsection{Меры советской власти в области культуры и в отношении Церкви за годы революции и Гражданской войны}

Несмотря на регулярное закрытие храмов и репрессии священнослужителей, с 1918 года Церковь также внесла помощь. Правительству это не нравится, и в 1922 году издаётся декрет, по которому происходит принудительное изъятие церковных ценностей.

\subsection{Исторические итоги Гражданской войны:}

\begin{enumerate}
    \item Огромные человеческие жертвы (около 13 млн. чел.).
    \item Величайшее хозяйственное разорение.
    \item Голод 1921--1923 г.
    \item Огромное количество беспризорников (8--9 млн. чел.). 
    \item Сохранение диктатуры Российской коммунистической партии большевиков.
\end{enumerate}

\subsection{Причины победы большевиков:}

\begin{enumerate}
    \item Энергичные действия советского правительства по созданию нового государства и армии.
    \item Продовольственная диктатура.
    \item Единство и централизация армии.
    \item Активная пропагандистская и идеологическая работа.
    \item Безжалостный террор.
\end{enumerate}

\subsection{Причины поражения белых:}

\begin{enumerate}
    \item Отсутствие единства в армии.
    \item Поддержка интервентов, расценивавшаяся как предательство.
\end{enumerate}
